\documentclass[12pt]{article}
\usepackage[utf8]{inputenc}
\usepackage[T1]{fontenc}
\usepackage[margin=1in]{geometry}
\usepackage{setspace}
\usepackage{xcolor}
\usepackage[colorlinks=true, urlcolor=blue!60!black, linkcolor=blue!60!black]{hyperref}
\setlength{\parindent}{2em}
\setlength{\parskip}{1.25em}
\renewcommand{\baselinestretch}{1.0}
\begin{document}
\begin{center}
{\LARGE \textit{New York State Rifle \& Pistol Association, Inc. v. Bruen}, 597 U.S. 1 (2022)\par}
\vspace{0.8em}
{\large Case Brief\par}
\end{center}
\vspace{1em}

\section{Memory Jogger}

The Second and Fourteenth Amendments protect an ordinary, law-abiding citizen's right to carry a handgun in public for self-defense, and firearm regulations covering protected conduct are constitutional only if the government proves they are consistent with the Nation's historical tradition of firearm regulation rather than under means-end scrutiny (\href{https://www.courtlistener.com/opinion/6480696/new-york-state-rifle-pistol-assn-inc-v-bruen/}{CourtListener}).

\section{Detailed Case Facts}

New York has regulated the public carriage of handguns since the early twentieth century, beginning with a 1905 misdemeanor prohibition on unlicensed concealed carry in cities and villages, followed by the 1911 ``Sullivan Law'' (1911 N.Y. Laws ch. 195, \S 1) criminalizing unlicensed handgun possession generally, and a 1913 amendment (1913 N.Y. Laws ch. 608, \S 1) requiring applicants for concealed-carry licenses to demonstrate ``good moral character'' and ``proper cause'' (\href{https://www.law.cornell.edu/supremecourt/text/20-843}{Cornell LII}). Under the modern codification at issue, N.Y. Penal Law Ann. \S 400.00(2)(f), possessing any firearm without a license was a crime, and because New York prohibited open carry of handguns altogether, obtaining an unrestricted concealed-carry license under \S 400.00(2)(f) was the sole lawful way for most residents to carry a handgun in public for self-defense (\href{https://www.law.cornell.edu/supremecourt/text/20-843}{Cornell LII}). Although no New York statute defined ``proper cause,'' New York courts interpreted the phrase to require an applicant to ``demonstrate a special need for self-protection distinguishable from that of the general community'' (\textit{In re Klenosky}, 75 App. Div. 2d 793 (1980)), which excluded generalized safety concerns or living in a high-crime neighborhood (\href{https://www.law.cornell.edu/supremecourt/text/20-843}{Cornell LII}). At the time of the decision, 43 states operated objective ``shall-issue'' licensing regimes, whereas only six states (New York, California, Hawaii, Maryland, Massachusetts, and New Jersey) and the District of Columbia maintained discretionary ``may-issue'' regimes (\href{https://www.law.cornell.edu/supremecourt/text/20-843}{Cornell LII}).

Petitioners Brandon Koch (of Troy, New York) and Robert Nash (of Averill Park, New York) were law-abiding adult residents of Rensselaer County and members of petitioner New York State Rifle \& Pistol Association, Inc. (\href{https://www.law.cornell.edu/supremecourt/text/20-843}{Cornell LII}). In 2014, Nash applied for an unrestricted public-carry license for general self-defense; in early 2015, the State granted him only a restricted license for hunting and target shooting, and in late 2016 a Rensselaer County licensing officer (New York Supreme Court Justice Richard McNally) denied Nash's request to lift those restrictions even after Nash cited recent neighborhood robberies (\href{https://www.law.cornell.edu/supremecourt/text/20-843}{Cornell LII}). Similarly, between 2008 and 2017, Koch held a restricted license for hunting and target shooting; when he applied in late 2017 to remove the restrictions based on his extensive safe firearms-handling experience, Justice McNally denied an unrestricted license and expanded Koch's permit only to allow carrying to and from work (\href{https://www.law.cornell.edu/supremecourt/text/20-843}{Cornell LII}).

\section{Procedural History}

Petitioners Koch, Nash, and the New York State Rifle \& Pistol Association sued state officials---originally Superintendent of the New York State Police George P. Beach II (later succeeded by Keith M. Corlett and then Kevin P. Bruen) and Justice Richard McNally---in the U.S. District Court for the Northern District of New York under 42 U.S.C. \S 1983 for declaratory and injunctive relief, alleging that the denial of unrestricted carry licenses for lack of ``proper cause'' violated the Second and Fourteenth Amendments (\href{https://www.law.cornell.edu/supremecourt/text/20-843}{Cornell LII}; \href{https://www.law.cornell.edu/supct/cert/20-843}{Cornell LII Preview}). Relying on binding Second Circuit precedent in \textit{Kachalsky v. County of Westchester}, 701 F.3d 81 (2d Cir. 2012), which had upheld New York's proper-cause requirement under intermediate scrutiny, the District Court dismissed the complaint (\textit{N.Y. State Rifle \& Pistol Ass'n v. Beach}, 354 F. Supp. 3d 143 (N.D.N.Y. 2018)) (\href{https://scholarship.law.duke.edu/cgi/viewcontent.cgi?article=1212&context=djclpp_sidebar}{Duke Journal of Constitutional Law \& Public Policy Sidebar}).

On appeal, the U.S. Court of Appeals for the Second Circuit affirmed by summary order, adhering to \textit{Kachalsky} (\textit{N.Y. State Rifle \& Pistol Ass'n v. Beach}, 818 F. App'x 99, 100 (2d Cir. 2020)) (\href{https://www.law.cornell.edu/supremecourt/text/20-843}{Cornell LII}). The U.S. Supreme Court granted certiorari on April 26, 2021 (No. 20-843), heard oral argument on November 3, 2021, and issued its decision on June 23, 2022, reversing and remanding the Second Circuit's judgment (\href{https://www.law.cornell.edu/supremecourt/text/20-843}{Cornell LII}; \href{https://www.supremecourt.gov/opinions/21pdf/597us1r54_7648.pdf}{Supreme Court Page Proof}).

\section{Judicial Votes}

The Court decided the case by a \textbf{6--3} vote (\href{https://www.law.cornell.edu/supremecourt/text/20-843}{Cornell LII}):
\begin{itemize}
    \item \textbf{Majority (6):} Justice Thomas delivered the opinion of the Court, joined in full by Chief Justice Roberts and Justices Alito, Gorsuch, Kavanaugh, and Barrett.
    \item \textbf{Concurrences:} Justice Alito filed a concurring opinion; Justice Kavanaugh filed a concurring opinion, joined by Chief Justice Roberts; Justice Barrett filed a concurring opinion.
    \item \textbf{Dissent (3):} Justice Breyer filed a dissenting opinion, joined by Justices Sotomayor and Kagan.
\end{itemize}

\section{Holding}

\textbf{Judgment:} Reversed and remanded (\href{https://www.law.cornell.edu/supremecourt/text/20-843}{Cornell LII}). New York's ``proper-cause'' requirement in N.Y. Penal Law Ann. \S 400.00(2)(f) violates the Fourteenth Amendment by preventing law-abiding citizens with ordinary self-defense needs from exercising their Second Amendment right to keep and bear arms in public (\href{https://www.law.cornell.edu/supremecourt/text/20-843}{Cornell LII}).

\textbf{Rule Announced:} Consistent with \textit{District of Columbia v. Heller}, 554 U.S. 570 (2008), and \textit{McDonald v. Chicago}, 561 U.S. 742 (2010), the Second and Fourteenth Amendments protect an individual's right to carry a handgun for self-defense outside the home (\href{https://www.law.cornell.edu/supremecourt/text/20-843}{Cornell LII}). Rejecting the two-step means-end scrutiny framework adopted by the federal courts of appeals, the Court announced a text-and-history standard: when the Second Amendment's plain text covers an individual's conduct, the Constitution presumptively protects that conduct, and the government bears the burden of justifying its regulation by demonstrating that it is consistent with the Nation's historical tradition of firearm regulation (\href{https://www.law.cornell.edu/supremecourt/text/20-843}{Cornell LII}). For modern regulations addressing unprecedented societal concerns or dramatic technological changes, courts use analogical reasoning asking whether modern and historical regulations are ``relevantly similar,'' guided by at least two central metrics: ``how and why'' the regulations burden a law-abiding citizen's right to armed self-defense (\href{https://www.law.cornell.edu/supremecourt/text/20-843}{Cornell LII}).

\textbf{Holding vs. Dicta:} The core holdings are that (1) the Second and Fourteenth Amendments guarantee an individual right to carry handguns in public for self-defense; (2) means-end scrutiny (such as intermediate or strict scrutiny) is forbidden in Second Amendment adjudication in favor of the text-and-history test; and (3) New York's discretionary ``proper-cause'' requirement is unconstitutional (\href{https://www.law.cornell.edu/supremecourt/text/20-843}{Cornell LII}). By contrast, Footnote 9's assurance that objective ``shall-issue'' licensing regimes (requiring background checks or firearms safety training) are generally permissible unless ``put toward abusive ends'' through lengthy wait times or exorbitant fees, the Court's observations about ``sensitive places'' (such as legislative assemblies, polling places, and courthouses, while rejecting declaring all of Manhattan a sensitive place), and the Court's discussion leaving open whether 1791 or 1868 controls incorporated rights are non-dispositive dicta or reservations not necessary to invalidate New York's statute (\href{https://www.law.cornell.edu/supremecourt/text/20-843}{Cornell LII}; \href{https://www.law.cornell.edu/constitution-conan/amendment-2/the-bruen-decision-and-concealed-carry-licenses}{Constitution Annotated}).

\section{Analysis of Opinions}

\subsection{Majority Opinion (Justice Thomas)}

Justice Thomas began by dismantling the two-step framework that federal courts of appeals had developed after \textit{Heller}, declaring it ``one step too many'' (\href{https://www.law.cornell.edu/supremecourt/text/20-843}{Cornell LII}). While Step One (determining whether regulated conduct falls within the Amendment's text as informed by history) aligned with \textit{Heller}, Step Two (applying strict or intermediate means-end scrutiny) contradicted \textit{Heller} and \textit{McDonald}, which rejected judge-empowering interest balancing because the Second Amendment itself is ``the very product of an interest balancing by the people'' (\href{https://www.law.cornell.edu/supremecourt/text/20-843}{Cornell LII}). Comparing the Second Amendment to First Amendment categorical exclusions and Sixth Amendment Confrontation Clause doctrine, Justice Thomas held that once the plain text covers an individual's conduct, the government must affirmatively prove that its regulation fits within this Nation's historical tradition of firearm regulation (\href{https://www.law.cornell.edu/supremecourt/text/20-843}{Cornell LII}). Where a regulation addresses a persistent societal problem (like urban handgun violence), the absence of a distinctly similar historical regulation is relevant evidence of unconstitutionality; where modern regulations address unprecedented societal concerns or technological changes, courts must reason by analogy to identify a ``well-established and representative historical analogue, not a historical twin,'' evaluating ``how and why'' the laws burden armed self-defense (\href{https://www.law.cornell.edu/supremecourt/text/20-843}{Cornell LII}).

Applying this framework, Justice Thomas first concluded that the plain text presumptively protected Koch's and Nash's conduct: both were ordinary, law-abiding members of ``the people,'' handguns are weapons ``in common use'' for self-defense, and the word ``bear'' naturally encompasses carrying arms outside the home where confrontation can occur (\href{https://www.law.cornell.edu/supremecourt/text/20-843}{Cornell LII}). Turning to history, Justice Thomas surveyed five eras invoked by New York and concluded that respondents failed to carry their burden: (1) medieval and early modern English law, including the 1328 Statute of Northampton, was either obsolete by the Founding or, after \textit{Sir John Knight's Case} (1686) and the 1689 English Bill of Rights, barred only carrying arms with evil intent to terrorize the public; (2) colonial and early Republic statutes (such as Massachusetts, New Hampshire, and East New Jersey) targeted terrorizing conduct or unusual weapons rather than peaceable carry of common self-defense arms; (3) antebellum common-law affray offenses, concealed-carry bans (which state courts upheld only when open carry remained available, as in \textit{Nunn v. State} (1846)), and mid-nineteenth-century surety statutes (which presumed a right to carry unless someone showed reasonable cause to fear a breach of the peace) never barred ordinary citizens from public carry; (4) Reconstruction-era discourse, including the Freedmen's Bureau Act, confirmed a right to public carry for self-defense, rendering Texas's 1871 ``reasonable grounds'' statute (\textit{English v. State} and \textit{State v. Duke}) an outlier; and (5) late-nineteenth-century Western territorial bans governed less than 1\% of the national population, were transitory, and could not override earlier constitutional understanding (\href{https://www.law.cornell.edu/supremecourt/text/20-843}{Cornell LII}). Declining to resolve whether 1791 or 1868 is the primary benchmark because the public understanding was the same in both periods, the Court invalidated New York's proper-cause requirement (\href{https://www.law.cornell.edu/supremecourt/text/20-843}{Cornell LII}).

\subsection{Concurring Opinion (Justice Alito)}

Justice Alito joined the Court's opinion in full and wrote separately to rebut Justice Breyer's dissent (\href{https://www.law.cornell.edu/supremecourt/text/20-843}{Cornell LII}). He emphasized the narrow scope of the holding: \textit{Bruen} decided only that a state may not enforce a discretionary proper-cause law that effectively bars law-abiding residents from carrying a handgun outside the home for self-defense, and it did not disturb \textit{Heller} or \textit{McDonald} regarding who may possess firearms, what weapons may be possessed, or purchase requirements (\href{https://www.law.cornell.edu/supremecourt/text/20-843}{Cornell LII}). Challenging the relevance of the dissent's statistics on mass shootings (including Buffalo), suicides, domestic disputes, and youth gun deaths, Justice Alito argued that New York's law did not deter criminals bent on violence while leaving ordinary citizens vulnerable, citing defensive gun use surveys and amicus briefs from vulnerable communities (\href{https://www.law.cornell.edu/supremecourt/text/20-843}{Cornell LII}). He also criticized means-end scrutiny as toothless, pointing to the Second Circuit's prior validation of New York City's premises-transport ban in \textit{NYSRPA v. City of New York} and Justice Breyer's \textit{Heller} dissent (\href{https://www.law.cornell.edu/supremecourt/text/20-843}{Cornell LII}).

\subsection{Concurring Opinion (Justice Kavanaugh, joined by Chief Justice Roberts)}

Justice Kavanaugh, joined by Chief Justice Roberts, joined the majority opinion while writing separately to underscore two limits on the decision (\href{https://www.law.cornell.edu/supremecourt/text/20-843}{Cornell LII}). First, \textit{Bruen} does not prohibit states from requiring licenses to carry handguns in public and leaves intact the objective ``shall-issue'' licensing regimes of 43 states, which may require fingerprinting, background checks, mental health records checks, and firearms safety and use-of-force training, subject only to as-applied challenges if abused (\href{https://www.law.cornell.edu/supremecourt/text/20-843}{Cornell LII}). Second, quoting \textit{Heller} and \textit{McDonald}, Justice Kavanaugh reiterated that the Second Amendment is ``neither a regulatory straightjacket nor a regulatory blank check'' and does not cast doubt on longstanding prohibitions on firearm possession by felons and the mentally ill, sensitive-place restrictions, conditions on commercial sales, or bans on ``dangerous and unusual weapons'' (\href{https://www.law.cornell.edu/supremecourt/text/20-843}{Cornell LII}).

\subsection{Concurring Opinion (Justice Barrett)}

Justice Barrett joined the Court's opinion in full and wrote separately to flag two unresolved methodological questions in originalist interpretation (\href{https://www.law.cornell.edu/supremecourt/text/20-843}{Cornell LII}). First, the Court did not settle when and how post-ratification practice (through liquidation, tradition, or precedent) may illuminate original public meaning, including how early such practice must begin and whether it applies to individual rights as well as structural provisions (\href{https://www.law.cornell.edu/supremecourt/text/20-843}{Cornell LII}). Second, she highlighted the unresolved scholarly debate over whether 1791 (ratification of the Second Amendment) or 1868 (ratification of the Fourteenth Amendment) supplies the controlling historical baseline for incorporated rights, cautioning against ``freewheeling reliance on historical practice from the mid-to-late 19th century'' (\href{https://www.law.cornell.edu/supremecourt/text/20-843}{Cornell LII}).

\subsection{Dissenting Opinion (Justice Breyer, joined by Justices Sotomayor and Kagan)}

Justice Breyer, joined by Justices Sotomayor and Kagan, opened with extensive empirical data on American gun violence---noting 45,222 firearm deaths in 2020, 277 mass shootings in the first half of 2022, firearm deaths surpassing motor vehicle crashes as the leading cause of death for children and adolescents, and heightened risks in urban areas, road rage, armed protests, domestic violence, and police encounters---to argue that elected legislatures must have constitutional leeway to address gun violence (\href{https://www.law.cornell.edu/supremecourt/text/20-843}{Cornell LII}). He identified three core errors in the majority opinion: (1) deciding the case on the pleadings without discovery or an evidentiary record showing how New York's proper-cause standard operated in practice; (2) discarding the unanimous two-step means-end framework of the courts of appeals in favor of a ``rigid history-only approach'' that misread \textit{Heller} and treated the Second Amendment differently from other constitutional rights; and (3) demonstrating the unworkability of the history-only test, which forces non-historian judges to resolve contested historical disputes and allows them to ``pick their friends out of history's crowd'' by discounting centuries of English, colonial, antebellum, surety, and nineteenth-century public-carry restrictions (\href{https://www.law.cornell.edu/supremecourt/text/20-843}{Cornell LII}; \href{https://www.law.cornell.edu/constitution-conan/amendment-2/the-bruen-decision-and-concealed-carry-licenses}{Constitution Annotated}).

\section{Examples: Future Applications}

\begin{enumerate}
    \item \textbf{Same-Side (Unconstitutional under \textit{Bruen}): Discretionary ``Good Moral Character'' Catch-All Licensing.} A state repeals its ``proper-cause'' requirement after \textit{Bruen} but enacts a statute empowering local police chiefs to deny a concealed-carry permit whenever an applicant lacks ``good moral character,'' defined as the ``essential character, temperament, and judgment necessary to be entrusted with a weapon,'' without objective disqualifying criteria. Because carrying a handgun in public is covered by the Second Amendment's plain text and this open-ended subjective standard replicates the unchanneled official discretion invalidated in \textit{Bruen} (rather than the ``narrow, objective, and definite standards'' described in Footnote 9), a court would strike down the discretionary catch-all as inconsistent with historical tradition (\href{https://www.law.cornell.edu/supremecourt/text/20-843}{Cornell LII}).

    \item \textbf{Same-Side (Unconstitutional under \textit{Bruen}): Blanket ``Sensitive Place'' Ban Covering Entire Commercial Districts.} A state designates ``all privately owned commercial establishments open to the public, public sidewalks in downtown business districts, and public transit hubs'' as gun-free ``sensitive places.'' Just as \textit{Bruen} rejected New York's attempt to treat the island of Manhattan as a sensitive place merely because people congregate there and police are present, this sweeping geographic ban eviscerates the general right to public carry and lacks a comparable historical analogue to narrow sensitive locations like courthouses, polling places, or legislative assemblies (\href{https://www.law.cornell.edu/supremecourt/text/20-843}{Cornell LII}).

    \item \textbf{Opposite-Side (Constitutional under \textit{Bruen}): Objective Shall-Issue Licensing with Fingerprinting, Background Check, and Safety Course.} A state requires every concealed-carry applicant to be at least 21 years old, pass a criminal and mental-health background check, submit fingerprints, pay a modest administrative fee, and complete an eight-hour certified firearms safety course, after which the state \textit{must} issue the permit within 30 days. Under \textit{Bruen}'s Footnote 9 and Justice Kavanaugh's concurrence, such objective shall-issue requirements ensure only that permit holders are ``law-abiding, responsible citizens'' and do not require any showing of special need, making the law presumptively constitutional absent proof of abusive delays or prohibitive fees (\href{https://www.law.cornell.edu/supremecourt/text/20-843}{Cornell LII}).

    \item \textbf{Opposite-Side (Constitutional under \textit{Bruen}): Temporary Disarmament of Individuals Judicially Found to Pose a Credible Physical Threat.} Congress prohibits firearm possession by a person subject to a domestic-violence restraining order that includes an individualized judicial finding that the person represents a credible threat to the physical safety of an intimate partner or child. As the Supreme Court later confirmed in \textit{United States v. Rahimi}, 602 U.S. 680 (2024), applying \textit{Bruen}'s analogical framework, this regulation is ``relevantly similar'' in ``how and why'' it burdens armed self-defense to Founding-era surety laws and ``going armed'' laws, which reflect a historical principle that individuals who threaten physical harm to others may be temporarily disarmed (\href{https://www.law.cornell.edu/supremecourt/text/22-915}{Cornell LII (Rahimi)}).

    \item \textbf{Fence-Sitter (Close / Unresolved under \textit{Bruen}): State Prohibition on Carrying Firearms on Private Property Open to the Public Unless the Owner Affirmatively Consents (``Vampire Rule'').} A state enacts a default rule making it a crime for a licensed handgun carrier to enter any privately owned business open to the public (such as a grocery store or gas station) unless the property owner posts a sign expressly permitting firearms. Proponents argue the law merely sets a state property-law default rule rooted in colonial and Reconstruction-era anti-poaching and enclosed-lands trespass statutes prohibiting carrying guns onto another's land without consent; challengers counter that those historical laws applied to private agricultural lands or plantations rather than commercial venues open to the general public, and that flipping the default effectively nullifies public carry across modern society. Because analogical similarity turns sharply on the chosen level of generality and whether 1791 or 1868 controls, this issue sits squarely on the constitutional fence (a question the Court subsequently took up in \textit{Wolford v. Lopez}) (\href{https://repository.uclawsf.edu/cgi/viewcontent.cgi?article=2238&context=hastings_constitutional_law_quaterly}{UC Law SF}; \href{https://www.law.cornell.edu/supremecourt/text/24-1046}{Cornell LII (Wolford)}).
\end{enumerate}

\section{Critique}

\textit{Bruen} sparked intense scholarly debate across ideological and methodological lines. From an \textbf{originalist and textualist perspective}, supporters such as Mark W. Smith praised \textit{Bruen} for faithfully enforcing \textit{Heller}'s text-and-history methodology and ending lower courts' use of deferential intermediate scrutiny to dilute an enumerated right (\href{https://journals.law.harvard.edu/jlpp/much-ado-about-nothing-rahimi-reinforces-bruen-and-heller-mark-w-smith/}{Harvard JLPP}). Yet even prominent originalist scholars raised sharp internal critiques: William Baude and Robert Leider argued that \textit{Bruen}'s rigid search for statutory analogues overlooks the Founding-era ``general law'' framework, which permitted reasonable regulation of the right to bear arms (\href{https://texaslawreview.org/applying-history-as-law-the-role-of-historical-facts-in-implementing-constitutional-doctrine/}{Texas Law Review}), while Randy Barnett, Lawrence Solum, and other originalists questioned the methodological tension between original public meaning and a free-standing ``history and tradition'' test, as well as \textit{Bruen}'s failure to resolve the 1791-versus-1868 incorporation baseline highlighted by Justice Barrett (\href{https://texaslawreview.org/applying-history-as-law-the-role-of-historical-facts-in-implementing-constitutional-doctrine/}{Texas Law Review}; \href{https://scholarship.law.wm.edu/cgi/viewcontent.cgi?article=4000&context=wmlr}{William \& Mary Law Review}).

From a \textbf{progressive and doctrinal perspective}, scholars such as Joseph Blocher, Eric Ruben, and Jacob D. Charles criticized \textit{Bruen}'s ``originalism-by-analogy'' for generating widespread confusion in lower courts over the proper level of generality, the significance of historical silence, and judicial competence to conduct ``law office history'' based solely on party-compiled records (\href{https://scholarship.law.duke.edu/cgi/viewcontent.cgi?article=6973&context=faculty_scholarship}{Duke Law Scholarship}; \href{https://firearmslaw.duke.edu/wp-content/uploads/2023/06/Bruen-and-the-Silent-Past-3.20.2023-Version.pdf}{Duke Center for Firearms Law}). Feminist and public-safety scholars further warned that tethering modern gun laws to eighteenth- and nineteenth-century legal traditions---eras when women's legal rights were minimal and domestic violence was largely unregulated---threatens vital modern protections for vulnerable populations (\href{https://www.wakeforestlawreview.com/wp-content/uploads/2024/06/59-Wake-Forest-L.-Rev.-131.pdf}{Wake Forest Law Review}). Indeed, lower courts' struggles with \textit{Bruen} prompted the Supreme Court two years later in \textit{United States v. Rahimi}, 602 U.S. 680 (2024), to clarify that \textit{Bruen} requires consistency with the ``principles that underpin our regulatory tradition'' rather than a ``historical twin'' (\href{https://www.law.cornell.edu/wex/second_amendment}{Cornell LII Wex}).

At bottom, \textit{Bruen} reflects a deep clash of \textbf{competing constitutional values}: protecting an individual's enumerated right to armed self-defense against discretionary state suppression and judicial balancing, versus preserving democratic responsiveness, federalism, and legislative flexibility to confront modern gun violence with empirical policy judgments.

\section{Key Quotations}

\begin{quote}
``When the Second Amendment's plain text covers an individual's conduct, the Constitution presumptively protects that conduct. The government must then justify its regulation by demonstrating that it is consistent with the Nation's historical tradition of firearm regulation. Only then may a court conclude that the individual's conduct falls outside the Second Amendment's `unqualified command.' ''
\end{quote}
\noindent --- Justice Thomas, Opinion of the Court (\href{https://www.law.cornell.edu/supremecourt/text/20-843}{Cornell LII}).

\begin{quote}
``To be clear, analogical reasoning under the Second Amendment is neither a regulatory straightjacket nor a regulatory blank check. On the one hand, courts should not `uphold every modern law that remotely resembles a historical analogue,' because doing so `risk[s] endorsing outliers that our ancestors would never have accepted.' On the other hand, analogical reasoning requires only that the government identify a well-established and representative historical analogue, not a historical twin. So even if a modern-day regulation is not a dead ringer for historical precursors, it still may be analogous enough to pass constitutional muster.''
\end{quote}
\noindent --- Justice Thomas, Opinion of the Court (\href{https://www.law.cornell.edu/supremecourt/text/20-843}{Cornell LII}).

\begin{quote}
``The constitutional right to bear arms in public for self-defense is not `a second-class right, subject to an entirely different body of rules than the other Bill of Rights guarantees.' We know of no other constitutional right that an individual may exercise only after demonstrating to government officers some special need.''
\end{quote}
\noindent --- Justice Thomas, Opinion of the Court (\href{https://www.law.cornell.edu/supremecourt/text/20-843}{Cornell LII}).

\begin{quote}
``Going forward, therefore, the 43 States that employ objective shall-issue licensing regimes for carrying handguns for self-defense may continue to do so. Likewise, the 6 States including New York potentially affected by today's decision may continue to require licenses for carrying handguns for self-defense so long as those States employ objective licensing requirements like those used by the 43 shall-issue States.''
\end{quote}
\noindent --- Justice Kavanaugh, Concurring Opinion, joined by Chief Justice Roberts (\href{https://www.law.cornell.edu/supremecourt/text/20-843}{Cornell LII}).

\begin{quote}
``In my view, when courts interpret the Second Amendment, it is constitutionally proper, indeed often necessary, for them to consider the serious dangers and consequences of gun violence that lead States to regulate firearms.''
\end{quote}
\noindent --- Justice Breyer, Dissenting Opinion, joined by Justices Sotomayor and Kagan (\href{https://www.law.cornell.edu/supremecourt/text/20-843}{Cornell LII}).

\end{document}
